\begin{textsample}{NEG Dim 8 – human – Score -2.00 – sp/sp\_ef\_intro\_2\_p6.txt}  \label{ex:f8_neg_002}
No estágio do desenvolvimento cognitivo compreendido dos 6 aos 12 anos, a criança passa a desenvolver conceitos mais \textbf{elaborados} em relação a ela mesma, apresentando maior controle emocional. É nessa fase que os conflitos aparecem, e a escola tem fundamental importância para que a criança passe a ampliar esse controle e as interações sociais construindo sua identidade socialmente, aprendendo a \textbf{avaliar} e a fazer escolhas para sua vida.
Dessa forma, amplia -se a autonomia intelectual, compreensão das normas e interesses pela vida social, promovendo a interação com sistemas mais amplos.
Nesse estágio há também uma expectativa em relação à produtividade do estudante em contraponto com o sentimento de inferioridade ; e o não cuidado para esses comportamentos, abre espaço para a baixa autoestima.
Dessa forma, um currículo voltado para o desenvolvimento das competências socioemocionais pode promover atividades que oportunizem aos estudantes lidar com esses sentimentos e assim desenvolver as habilidades como a resiliência e a empatia.

% matched lemmas: avaliar, elaborar
\end{textsample}
